The \textsc{bisect} platform generates all six K-track compactness metrics
in a single command:
\begin{verbatim}
bisect label-analyze <label> --year 2020 --types all
\end{verbatim}

All metrics are implemented in \texttt{crates/bisect-analysis/src/compactness.rs}
and returned together via the \texttt{all\_metrics()} function:
\begin{itemize}
  \item \texttt{polsby\_popper()} — K.1, range $[0,1]$
  \item \texttt{reock()} — K.2, range $[0,1]$ (centroid approximation)
  \item \texttt{convex\_hull\_ratio()} — K.3, range $[0,1]$
  \item \texttt{schwartzberg()} — K.4, range $[1,\infty)$
  \item \texttt{length\_width\_ratio()} — K.5, range $[1,\infty)$
  \item \texttt{population\_weighted\_compactness()} — K.6, separate function (requires tract populations)
\end{itemize}
\texttt{reock()} is a production parity metric, not an exact polygon-MBC
implementation.  Expert reports should disclose the centroid-radius approximation
when comparing against software that reports exact-MBC Reock.
For exact-MBC reference evidence, use the smoke package at
\texttt{docs/examples/k-reock-evidence-packages/exact-mbc-smoke/}; it validates
the exact helper and shows where the production proxy can diverge from exact
polygon-MBC Reock.  It is not a real district-polygon replay package.

\subsection{Output Format}

The output is a JSON file at \texttt{analysis/<label>/<year>/compactness.json}:
\begin{verbatim}
{
  "label": "nc_ratio_optimal_2020",
  "year": "2020",
  "seed": 42,
  "projection": "EPSG:5070",
  "algorithm": "ratio-optimal",
  "districts": [
    {
      "district_id": 1,
      "population": 745234,
      "polsby_popper": 0.231,
      "reock": 0.374,
      "convex_hull": 0.952,
      "schwartzberg": 2.079,
      "length_width": 1.647,
      "pwc_km2": 58.3,
      "within_acceptable_range": true
    },
    ...
  ],
  "state_summary": {
    "mean_pp": 0.224,
    "mean_reock": 0.371,
    "mean_ch": 0.947,
    "mean_schwartzberg": 2.113,
    "mean_lw": 1.683,
    "mean_pwc_km2": 67.8,
    "max_lw": 2.43,
    "districts_exceeding_lw3": 0,
    "all_ch_above_085": true
  }
}
\end{verbatim}

\subsection{Verification Chain}

The JSON output is linked to the plan's SHA-256 hash chain:
\begin{verbatim}
bisect label-verify <label> --year 2020
\end{verbatim}
verifies that the compactness output corresponds to the exact district
assignment file and TIGER/Line data that produced it.
This provides a cryptographic chain of custody suitable for inclusion in expert reports.
Formal certification as a court-filing appendix requires counsel review and
expert witness attestation; the structured JSON output is the technical input
to that process, not the certified artifact itself.

\subsection{Acceptable Range Flags}

The ``acceptable range'' flags in the JSON are based on:
\begin{itemize}
  \item PP $\geq 0.18$ (conservative threshold from K.1)
  \item Reported Reock $\geq 0.30$ (from K.2)
  \item CH $\geq 0.82$ (from K.3)
  \item S $\leq 2.4$ (equivalent to PP $\geq 0.17$, from K.4)
  \item LW $\leq 2.8$ (court threshold $< 3$, from K.5)
  \item PWC within 2 standard deviations of state mean (from K.6)
\end{itemize}
These thresholds are conservative and can be adjusted for specific jurisdictions.
